\subsection{System Constraints and Limitations}

The development and operation of the Dikaios Amanah system are subject
to several technical, operational, security, and project-level
constraints. These constraints define the boundaries of the proposed
system and its implementation within the scope of the Software
Engineering Lab project.

\subsubsection{Project Scope Constraint}

Dikaios Amanah is developed as an academic Software Engineering Lab
project. Therefore, the system is intended to demonstrate the proposed
software architecture, functionality, security concepts, and workflow
rather than operate as a production financial platform.

Production-scale deployment, commercial operation, and regulatory
certification are outside the current project scope.

\subsubsection{Financial Transaction Limitation}

The investment functionality implemented within the project represents
simulated investment activities. The system does not process real
financial transactions, transfer real funds, or provide actual
investment services.

Any financial values displayed by the system are intended for
demonstration and testing purposes.

\subsubsection{Regulatory and Legal Constraints}

A production implementation of a platform involving investment
activities would be subject to applicable financial, business,
data-protection, and other regulatory requirements.

Obtaining licenses, regulatory approvals, legal certifications, and
formal authorization to operate a financial service are outside the
scope of this academic project.

\subsubsection{Third-Party Service Dependency}

The system may depend on external services for hosting, cloud
infrastructure, file storage, authentication, email or notification
delivery, and other supporting functionality.

Availability, pricing, service limits, and changes to third-party
services may affect system operation.

\subsubsection{Internet and Infrastructure Dependency}

The web-based system requires network connectivity for communication
between the client application, backend services, database, and
supporting cloud infrastructure.

Temporary network failures or service outages may prevent users from
accessing certain system functionality.

\subsubsection{Document Verification Limitation}

The system can collect and manage documents submitted for verification,
but the academic implementation cannot guarantee that submitted
documents are legally authentic.

Document review and verification functionality is intended to
demonstrate the software workflow and access-control mechanisms rather
than replace formal legal or governmental verification procedures.

\subsubsection{Security Constraints}

The system shall implement appropriate authentication, authorization,
input validation, secure credential handling, and audit mechanisms
within the project's technical scope.

However, the academic implementation cannot guarantee complete
protection against every possible cybersecurity threat. Production-level
penetration testing, security certification, continuous security
monitoring, and advanced threat detection are outside the current
scope.

\subsubsection{Scalability Constraint}

The initial implementation is designed for demonstration and academic
evaluation rather than large-scale commercial usage.

The system may therefore require additional infrastructure,
optimization, caching, database scaling, and load-balancing mechanisms
before supporting a large number of concurrent users.

\subsubsection{Data Availability Constraint}

The quality of system demonstrations and testing depends on the
availability of suitable test data.

The project may use controlled, synthetic, or sample data for
demonstration purposes. Real confidential business, financial, or
personal information shall not be required for the academic
implementation.

\subsubsection{Development Resource Constraint}

The project is developed within the time, hardware, software, and
resource limitations of an academic Software Engineering Lab project.

These limitations may restrict the number of advanced features that can
be implemented during the project period.

\subsubsection{Deployment Constraint}

The final system may be deployed using cloud-based infrastructure
appropriate for the project's requirements. However, the deployment
configuration is intended for demonstration and evaluation and may not
provide the redundancy, monitoring, disaster recovery, and operational
support expected from a production financial platform.